%Microservice architecture has gradually becomes a hot spot in corporate practice and academic research recently, many platforms, tools and theories have emerged to solve new problems faced by microservice systems.
%At present, there are several popular fault location methods for distributed systems, cloud and microservice systems, can be divided into the following two parts.

Traditional anomaly detection approaches for distributed and cloud-based systems are mainly based on log analysis.
Fu \emph{et. al}\cite{09ICDM_EADULA}  proposed an unstructured log analysis technique of anomalies detection for distributed systems. 
%It detects anomalies with the text converting algorithm, workflow model and performance measurement model.
Xu \emph{et al.}\cite{17ICEBE_LogDC} introduced LogDC, a log model based problem diagnosis tool for cloud applications with the full-lifecycle Kubernetes logs.
Jia \emph{et al.}\cite{17ICWS_ADHGML} proposed an approach for automatic anomaly detection based on logs. It raises anomaly alerts on observing deviations from the hybrid model which captures normal execution flows inter and intra services.
These approaches build anomaly detection modes with log parsing.
They cannot detect the anomalies with quality metrics of services, nor can they support anomaly propagation analysis for root cause localization.

Root cause analysis in practice often relies on visualization of logs and traces.
Zhou \emph{et al.}~\cite{Zhou_TSE18MS} presented an improved visualization analysis approach for fault analysis and debugging of microservice systems.
Wang \emph{et al.}~\cite{Hanzhang_GraphBased2019} demonstrated Grano, an end-to-end anomaly detection and root cause analysis system for cloud native distributed data platform by providing a holistic view of the system component topology, alarms and application events.
Guo \emph{et al.}~\cite{Guo_FSE20Industry} proposed a graph-based trace analysis approach for understanding architecture and diagnosing problems of microservice systems.
There approaches the developers to manually detect the root causes with the aid of trace and log visualization.

Recently, trace analysis based approaches have been proposed for automatic root cause localization for microservice or service-based systems.
%\textbf{Trace-based} 
%Trace-based approaches locating root causes by extracting features inside traces and analyzing the potential relevance of traces and anomalies.
Pham \emph{et al.}\cite{Cuong_traceBased} introduced a fault localization approach based on trace similarity comparison. It collects a large number of traces labelled with fault types and root causes through fault injection and conducts fault localization by comparing between faulty traces and historical traces.
Gan \emph{et al.}\cite{19ICASPLOS_Seer} presented Seer, an online cloud performance debugging system for QoS violations by learning spatial and temporal patterns from traces.
Zhou \emph{et al.}\cite{FSE2019LatentError} proposed MEPFL, an approach that can predict latent errors and locate root causes for microservice applications by learning from system trace logs.
It predicts latent errors, faulty microservices, and fault types for trace instances captured in the production environment using models trained based on a set of features defined on microservice traces.
These approaches often need to collect and analyze a large number of traces to extract anomaly patterns or train prediction models, thus are not efficient for large-scale microservice systems.

Some researchers have proposed approaches that use service call graphs or causality graphs for automatic root cause localization for microservice or service-based systems.
{\'{A}}lvaro \emph{et al.}~\cite{GraphBasedformicroservice} presented a root cause analysis framework, based on graph representations of service-oriented and microservice architectures. It analyzes anomalies by comparing the similarity between the fault subgraph and the historical subgraph.
Li \emph{et al.}~\cite{NOMS2020MicroRCA} proposed MicroRCA, a system to locate root causes of performance issues in microservice systems. MicroRCA infers the root causes by correlating the performance symptoms with the corresponding resource utilization based on an attributed graph.
Wang \emph{et al.}~\cite{DCCGRID2018CloudRange} proposed CloudRanger, a system dedicated for cloud native systems. CloudRanger identifies the culprit services which are responsible for cloud incidents by constructing the impact graph among services with a dynamic causal relationship analysis approach. 
Kim \emph{et al.}~\cite{cscs2013monitorRank} introduced MonitorRank, an algorithm that can reduce the time, domain knowledge, and human effort required to find the root causes of anomalies in service-oriented architectures. MonitorRank finds root causes with service call graph and historical and current time-series metrics of services.
Lin \emph{et al.}~\cite{ICSOC2018Microscope} proposed Microscope, a system to identify and locate the anomalous services with a ranked list of possible root causes in microservice environments. Microscope infers the causes of performance problems in real time by constructing a service causal graph.
The preceding approaches share some similar ideas with ours.
However, most of these approaches detect anomalous services by simply analyzing the correlations between metrics or statistical properties and do not consider the characteristics of different anomaly types and quality metrics. 
%Also these methods hardly consider the performance of anomaly detection.
Moreover, these approaches need to traverse most of the nodes in the graph, thus may be inefficient for large-scale systems.
In contrast, we combine machine learning and statistical methods and design a customized model for the detection of each type of anomalies and adopt a pruning strategy to improve the efficiency of anomaly propagation chain analysis.

